\begin{abstract}
This note responds to \emph{Approximation bridges} (relayed 10 October 2026). Its finite claims were re-checked
exactly; all pass. The bridges are then specialised to the bias-free ReLU network in three theorems. The first is an
exact identity: every layerwise scheme's final-mean error equals
$\sum_\ell\Gamma_\ell\,\Delta_\ell$. Here $\Delta_\ell$ is the scheme's error in the absolute moments $\E|z_\ell|$, and
$\Gamma_\ell=2^{-(L-\ell+1)}W_L\cdots W_{\ell+1}$ is a gate-free, state-free transport fixed by the weights. On the official
net 0 its operator norm equals the Fuss--Catalan edge $\tfrac12\,2^{-k/2}((k+1)^{k+1}/k^k)^{1/2}$ to within $6\%$ at every
depth $k=L-\ell$: $0.50,0.70,0.64,0.54,0.43,\dots,0.018$. The second identifies the final-only localisation's terminal
covariance actions for quadratic input features with the Hessian means $\E\nabla^2F$, which are wall sums, and proves a
transfer lemma: driving the deformation with approximate actions costs exactly their error integrated against the removed
interaction. The third gives the cost of exact local spherical actions on the network, about $2\sqrt2\,L$ wall crossings
per coordinate circle. Together these turn the open realisation theorem into one precise requirement: the late-layer laws,
projected on the next layer's rows, must be accurate in the kink (resolvent) test, weighted by $\|\Gamma_\ell\|$.
\end{abstract}

\maketitle

\section{The bridges' finite claims, checked}

Each claim below was re-derived and checked exactly, symbolically or with exact rationals; where a law is continuous, by
adaptive quadrature (\texttt{theory25/check\_bridges.py}). None of these is a network experiment.

\begin{center}\footnotesize
\begin{tabular}{@{}p{4.4cm}p{8.8cm}l@{}}\toprule
claim & check & result\\\midrule
double occupancy under a $\mathrm{Beta}(\frac12,\frac12)$ split & $ii:ij:jj=\frac38:\frac14:\frac38$; $\E\cos^4=\frac38$ & pass\\
complete-graph quartic mode & $H_{K_d}(S-\frac3{d+2})=\frac{d+2}4(S-\frac3{d+2})$, $d=3,4,6$, symbolic & pass\\
final-only drift & It\^o drift of $m_t$ equals $-V\T C^2a$; zero when $V\T C=0$ & pass\\
25-query resolvent quadrature & $M=12$: worst error $2.98\times10^{-5}\sigma$ over nine laws (atoms at $0$, $t_3$) $\le4.95\times10^{-5}\sigma$ & pass\\
variational bracket & $q_s\in[1-s^2J-\|r\|^2,\,1-s^2J]$, width $\le\|r\|^2$, 500 cases & pass\\
local leakage & $\|Pe(A)P\|\le\delta+(1-\delta)\min\{1,\eta^2/(a-\tau)^2\}$, 300 cases & pass\\
constant escape profile & $A\mathbf 1_E=a\mathbf 1_E\Rightarrow$ one step exact for every reward & pass\\\bottomrule
\end{tabular}
\end{center}

The It\^o computation is worth recording because it also gives the martingale version used in
Lemma~\ref{lem:transfer}. For $p\propto\exp(\frac12\phi\T J\phi+h\T\phi)\nu$ we have $\nabla_hm=V\T$,
$\nabla_Jm:X=\frac12\cov(F,\phi\T X\phi)$ and $\nabla_h^2m:X=\kappa_3(F,\phi,\phi):X$. Since
$\cov(F,\phi\T C^2\phi)=\kappa_3:C^2+2V\T C^2a$, the drift under $dJ=-C^2dt$, $dh=C\,dB$ is $-V\T C^2a$. Adding $C^2a\,dt$
to $dh$ removes it, so $dm=V\T C\,dB$.

\section{The final means are a fixed linear readout of absolute moments}

Write $z_\ell=W_\ell a_{\ell-1}$, $a_\ell=\relu(z_\ell)$, $a_0=X$ with $\E X=0$, $m_\ell=\E a_\ell$ and
$A_\ell=\E|z_\ell|$ (componentwise).

\begin{theorem}[Ungated recursion]\label{thm:ungated}
$m_\ell=\frac12(W_\ell m_{\ell-1}+A_\ell)$, hence
\[
m_L=\sum_{\ell=1}^L\Gamma_\ell A_\ell,\qquad \Gamma_\ell=2^{-(L-\ell+1)}\,W_L\cdots W_{\ell+1}\quad(\Gamma_L=\tfrac12I).
\]
\end{theorem}
\begin{proof}
$\relu(z)=\frac12(z+|z|)$; take expectations and unroll from $m_0=0$.
\end{proof}

\begin{theorem}[Exact error of any layerwise scheme]\label{thm:err}
Let a scheme carry model laws $\hat\nu_\ell$ of $a_\ell$ with means $\hat m_\ell$, and set $\hat m_\ell=\E_{\hat\nu_{\ell-1}}
\relu(W_\ell a)$. This covers every closure, mixture state and cumulant chain. With
$\Delta_\ell=\E_{\hat\nu_{\ell-1}}|W_\ell a|-\E_{\nu_{\ell-1}}|W_\ell a|$,
\[
\hat m_L-m_L=\sum_{\ell=1}^L\Gamma_\ell\,\Delta_\ell\qquad\text{exactly.}
\]
\end{theorem}
\begin{proof}
The scheme satisfies Theorem~\ref{thm:ungated} with $\hat A_\ell=\E_{\hat\nu_{\ell-1}}|W_\ell a|$, because its means are the
means of its laws. Subtract the two unrolled sums.
\end{proof}

Checked to $6\times10^{-17}$ for a Gaussian closure on a network with 2-D input, with the input law replaced by a
2-million-point angular quadrature; the identity is algebraic, so it holds for any law. The closure error itself is $0.15$
there.

\begin{proposition}[The transport contracts]\label{prop:gamma}
$\Gamma_\ell$ depends only on the weights. For i.i.d.\ $N(0,2/n)$ weights and $k=L-\ell$, the edge of the order-$k$
Fuss--Catalan law gives
\[
\|\Gamma_\ell\|_{\rm op}\to\gamma_k=\tfrac12\,2^{-k/2}\Big(\tfrac{(k+1)^{k+1}}{k^k}\Big)^{1/2},
\qquad \tfrac1{\sqrt n}\|\Gamma_\ell\|_F\to\tfrac12\,2^{-k/2}.
\]
Here $\gamma_k\le0.71$ for every $k$; $\gamma_0=\frac12$, $\gamma_1=0.707$, $\gamma_4=0.437$, $\gamma_8=0.150$ and
$\gamma_{15}=0.018$. On the official net 0 the exact norms, computed by power iteration on the weights, are
\[
\begin{aligned}
&0.500,\ 0.699,\ 0.637,\ 0.537,\ 0.429,\ 0.329,\ 0.249,\ 0.193,\\
&0.148,\ 0.107,\ 0.080,\ 0.060,\ 0.045,\ 0.033,\ 0.024,\ 0.0175
\end{aligned}
\]
for $k=0,\dots,15$, within $6\%$ of $\gamma_k$, and $\sum_k\|\Gamma\|\approx4.1$.
\end{proposition}
\begin{proof}
The squared singular values of a product of $k$ independent Ginibre matrices with entry variance $1/n$ converge to the
Fuss--Catalan law of order $k$, supported on $[0,(k+1)^{k+1}/k^k]$. Rescale by $2^{k/2}$ and multiply by $2^{-(k+1)}$. The
Frobenius statement is $\E\|Wv\|^2=2\|v\|^2$ iterated.
\end{proof}

No product of layer operator norms appears. The gated tangent $\prod D_\Phi W$ that carries mean errors in the closure
has gains near one at depth. The ungated transport is a contraction at every depth, and it is known before any state is
computed.

\begin{theorem}[Resolvent readout through depth]\label{thm:readout}
Let $\sigma_{\ell,i}^2\ge\E z_{\ell,i}^2$ and $q_{\ell,i,j}=\E\big[z_{\ell,i}^2/(z_{\ell,i}^2+\sigma_{\ell,i}^2e^{2jh})\big]
\in[0,1]$, with $h=\pi/\sqrt M$ and $|j|\le M$. Then
\[
m_L=\sum_{\ell=1}^L\Gamma_\ell\Big(\tfrac{2h}{\pi}\,\sigma_\ell\circ\sum_{j=-M}^Me^{jh}q_{\ell,\cdot,j}\Big)+R,\qquad
|R_i|\le\sum_\ell(|\Gamma_\ell|\,2\varepsilon_M\sigma_\ell)_i,
\]
with $\varepsilon_M=\big[\tfrac2{1-e^{-\pi\sqrt M}}+\tfrac2\pi\big]e^{-\pi\sqrt M}$ ($4.95\times10^{-5}$ at $M=12$). If every
query is bracketed, $q\in[q^-,q^+]$, the midpoint readout has error at most
$\sum_\ell|\Gamma_\ell|\,\frac{2h}{\pi}\,\sigma_\ell\circ\sum_je^{jh}\frac{q^+-q^-}2$ componentwise, plus $|R|$.
\end{theorem}
\begin{proof}
$|z|=\frac2\pi\int_0^\infty z^2/(z^2+t^2)\,dt$. Substitute $t=\sigma e^s$ and apply the companion's trapezoid bound to $\E|z|=2\E
z_+-\E z$, which doubles its constant. Then insert into Theorem~\ref{thm:ungated}.
\end{proof}

\begin{corollary}[A final-only certificate along actual rows]\label{cor:sliced}
\[
\operatorname{RMS}(\hat m_L-m_L)\le\sum_{\ell=1}^L\|\Gamma_\ell\|_{\rm op}\,\operatorname{RMS}_i\,
W_1\big(\mathcal L_{\hat\nu_{\ell-1}}(w_{\ell,i}\!\cdot a),\,\mathcal L_{\nu_{\ell-1}}(w_{\ell,i}\!\cdot a)\big).
\]
Only one-dimensional projected laws along the next layer's actual rows enter, tested by $|\cdot|$, the kink. Their weights
are $0.50, 0.70, 0.64, 0.54, 0.43$ for the last five layers and fall to $0.018$ at the first.
\end{corollary}
\begin{proof}
$|\cdot|$ is $1$-Lipschitz, so $|\Delta_{\ell,i}|$ is at most the $W_1$ distance of the projections. Apply
Theorem~\ref{thm:err} with the triangle inequality.
\end{proof}

\begin{remark}[Inheritance is relocated, not removed]
If $\hat\nu_{\ell-1}$ differs from $\nu_{\ell-1}$ only by a mean shift $\delta m$, then
$\Delta_\ell\approx\operatorname{diag}(2\Phi(\alpha_\ell)-1)\,W_\ell\,\delta m$. The identity of Theorem~\ref{thm:err} then
gives $\delta m_\ell=\operatorname{diag}\Phi(\alpha_\ell)W_\ell\,\delta m$, the gated tangent. The theorems therefore do not
say that early errors vanish. They say where the final readout spends accuracy: in the late-layer projected laws, at the
kink, with known weights. Errors born early count only through the late-layer projected laws they distort.
\end{remark}

\section{Final-only deformation, specialised}

\begin{proposition}[Terminal actions for quadratic features are Hessian means]\label{prop:price}
For $X\sim N(0,I_d)$ and $\phi=\operatorname{svec}(XX\T-I)$, Stein's identity applied twice gives
$\cov(\phi,F_i)=\operatorname{svec}\E\nabla^2F_i$. For the ReLU network,
\[
\E\nabla^2F_i=\sum_{\ell,j}\E\Big[\delta(z_{\ell,j})\,(A_\ell e_j)_i\,\nabla z_{\ell,j}\nabla z_{\ell,j}\T\Big],
\qquad \tr\E\nabla^2F_i=\mu_i,
\]
with $A_\ell$ the downstream gated product. These are the wall currents of the companion's eq.~(10), with the normal
outer product in place of its trace. Equivalently, $\partial\mu/\partial\Sigma=\frac12\E\nabla^2F$ at $\Sigma=I$ (Price).
The invisible quadratic deformations are the symmetric $\Delta$ with $\tr(\Delta\,\E\nabla^2F_i)=0$ for all $i$: codimension
at most $n$ among $d(d+1)/2$. A closure supplies approximate actions through its adjoint.
\end{proposition}

\begin{lemma}[Transfer, not removal]\label{lem:transfer}
Run the martingale version ($dJ=-C^2dt$, $dh=C\,dB+C^2a\,dt$) with $C_t\perp\hat V_t$, where $\hat V_t$ are approximate
actions. Then exactly
\[
\E\|m_T-m_0\|^2=\E\int_0^T\big\|(V_t-\hat V_t)\T C_t\big\|_F^2\,dt\;\le\;\sup_t\|V_t-\hat V_t\|_{\rm op}^2\;\E\tr(J_0-J_T).
\]
\end{lemma}
\begin{proof}
$dm=V\T C\,dB=(V-\hat V)\T C\,dB$, then It\^o's isometry. Also $\|E\T C\|_F\le\|E\|_{\rm op}\|C\|_F$, and
$\int\tr C^2\,dt=\tr(J_0-J_T)$.
\end{proof}

Consequences for this network. The Gaussian input carries no interaction, $J_0=0$, so any final-only deformation
starts from an interaction we introduce. That is the companion's gauge trap. Lemma~\ref{lem:transfer} prices every such
path: the drift is at most the action error (operator norm) times the square root of the interaction removed, and for
errors without special alignment it is of that order. An endpoint that differs materially from the start requires removing a material interaction. The
actions must then be accurate to about the target RMS divided by that amount. The mechanism therefore moves the accuracy
requirement from the means to the terminal actions. It removes the requirement only where the actions are exact, for
instance a single layer, where $\mu_i=\sqrt{w_i\T\Sigma w_i}/\sqrt{2\pi}$ and its Price derivative are closed-form, so the
constraints are explicit. At depth two and beyond no exact action is available.

\section{Exact local actions and what they cost}

\begin{proposition}[Coordinate circles cross few walls]
For uniform $\theta$ and a uniform coordinate pair $(i,j)$, the circle obtained by rotating $(\theta_i,\theta_j)$ crosses
in expectation approximately $2\pi\cdot\frac{\sqrt2}d\sum_{\rm gates}r_g$ walls of the whole network, $r_g$ being each gate's rotation
flip rate. With $r_g\approx1/\pi$, which holds exactly at layer 1 and in aggregate at every depth because $f'(1)=1$, this
is $2\sqrt2\,L\,n/d\approx45$ for $n=d=1024$, $L=16$. An exact action $Q_{ij}F(\theta)$ by circle traversal therefore costs
about $25$ forward passes. Each crossing is a rank-one update of the network's two-column restriction.
\end{proposition}
\begin{proof}
Rice's formula along the circle. The velocity has length $r_{ij}=\sqrt{\theta_i^2+\theta_j^2}$, with $\E r_{ij}\approx
\sqrt{2/d}$, and lies in a random coordinate plane. Its projection on a generic gradient has mean absolute value
$\|\nabla z\|\sqrt{2/(\pi d)}$, against $\|\nabla z\|\sqrt{2/\pi}$ for the full rotation, so the ratio is $\sqrt2/d$.
\end{proof}

The certified angular gap of the collision process is therefore real, and an exact local action on the terminal
observable costs about $25$ forward passes, not one. A Markov-chain use pays this per step and mixes at a rate of order
$1/d$ per collision. The process controls relaxation; it does not by itself evaluate $\E F$ more cheaply than sampling. The
exchange element $\Omega_G$ of the sparsifier bridge acts through relations that the network's layer maps break: no
positive map sends the transpositions to gate products, by the companion's three-slot obstruction. The leakage and clock
bounds apply as soon as a reversible representation with cheap actions exists. None of the exact primitives below
supplies one at depth.

\section{The remaining requirement, precisely}

Three primitives have exact and inexpensive actions on this problem:
\begin{itemize}[leftmargin=1.5em]
\item Gaussian integrals at the input;
\item linear algebra in the weights, including $\Gamma_\ell$;
\item network evaluations, forward passes and Jacobian--vector products, at chosen inputs.
\end{itemize}
Any certified computation of the final means is assembled from these. Theorem~\ref{thm:err} and
Corollary~\ref{cor:sliced} state the requirement exactly: produce model laws whose \emph{one-dimensional projections along
the next layer's rows} are accurate in the kink test, with errors $\varepsilon_\ell$ satisfying
$\sum_\ell\|\Gamma_\ell\|\varepsilon_\ell\le$ target. The weights concentrate the demand on the last five layers, where
$\|\Gamma\|\in[0.43,0.70]$. Earlier layers matter only through the late-layer projections they shape.

This fixes the joint construction the bridges ask for.
\begin{itemize}[leftmargin=1.5em]
\item \textbf{State.} The projected laws $\mathcal L(w_{\ell,i}\cdot a_{\ell-1})$, $n$ per layer, together with the shared
structure needed to transport them: covariance, the collective coordinate, and the incoherent per-neuron shape.
\item \textbf{Moves.} The layer map acting on that state. A projection of the next pre-activation is the dependent sum
$\sum_kw_k\relu(z_k)$, so its law needs pair structure (Price series), the collective mixture, and diagonal cumulant
contractions $\sum_kw_k^3\kappa_3(\relu z_k)$.
\item \textbf{Solver.} The $25$ bounded resolvent queries per neuron and layer of Theorem~\ref{thm:readout}, combined by
$\Gamma$.
\item \textbf{Certificate.} Projected-discrepancy bounds per layer weighted by $\|\Gamma_\ell\|$. Their widths direct
refinement to the neurons and layers whose weighted brackets dominate.
\end{itemize}
The decisive theorem is then a quantitative statement about one object. It is a bound, at the kink scale and conditional
on the collective structure, on the error of an affordable state's projected laws for dependent sums
$\sum_kw_k\relu(z_k)$ at the last five layers: a quantitative central limit theorem for these sums at resolution about
$10^{-4}$.

One measured fact already points the second computation at the right target. On net 0 the oblique (whitened) oracle
calibration of CC1's error reaches $8\times10^{-9}<3.1\times10^{-8}$ at the floor budget, with $82\%$ of its
signal-to-noise beyond the top $256$ output modes, which hold only $14\%$ of its error energy. A template for the quiet component of the first pass's error,
for instance the $\Gamma$-transported late-layer kink corrections of Theorem~\ref{thm:readout}, would need only one
whitened Monte Carlo coefficient.
